\documentclass[11pt,a4paper]{article}
\usepackage[margin=24mm]{geometry}
\usepackage{amsmath,amssymb,booktabs}
\usepackage[hidelinks]{hyperref}
\usepackage{microtype}
\newcommand{\dd}{\mathrm d}
\newcommand{\ii}{\mathrm i}
\newcommand{\cM}{\mathcal M}
\newcommand{\cL}{\mathcal L}
\title{Route G2: Local Scattering with a Recoiling Quantum Detector}
\author{Bounded tree-level effective-theory result}
\date{23 September 2026}
\begin{document}
\maketitle
\begin{abstract}
A positive local interaction between the G1 field and a dynamical
massive detector field permits a real change of Kaluza--Klein mode.
The detector is localized by its quantum state, not pinned by an
external potential. It carries the compensating compact momentum
and ordinary recoil. We derive the vertex, threshold, cross section,
wavepacket selection rule and energy balance. One engineering example
has a nonzero conversion cross section. This is a collision of two
systems, not coordinate-induced mass change or an elementary-particle
identification. No external drive or mass fit is used.
\end{abstract}

\section{One local action and its stability scope}
Keep the G1 background $\mathbb M^{3,1}\times S^1$ with radius $R$,
signature $(+----)$, $\theta\sim\theta+2\pi$ and
$D_\theta\Phi=(\partial_\theta+\ii\alpha)\Phi$.
Switch off its diagnostic external sources during scattering. Add a
complex field $X$ neutral under that connection:
\begin{align}
 S={}&\int\dd^4x\,R\dd\theta\,
 \bigg[
 |\partial_\mu\Phi|^2-\frac{|D_\theta\Phi|^2}{R^2}
 -M_5^2|\Phi|^2
 +|\partial_\mu X|^2-\frac{|\partial_\theta X|^2}{R^2}
 -M_D^2|X|^2-\kappa_5|\Phi|^2|X|^2
 \bigg], \label{eq:action}\\
 &M_5>0,\qquad M_D>0,\qquad \kappa_5\geq0 .
\end{align}
The interaction is point-local in all five coordinates. A wavepacket
of $X$ will localize the apparatus around the circle. This is
\emph{state localization}, not a delta-function defect in the action,
a prescribed classical background, or an infinitely heavy momentum
sink. The $X$ field has its own global $U(1)$ charge; net charge, not
total particle plus antiparticle number, is exactly conserved.

Both kinetic terms have positive sign. The classical potential
\begin{equation}
 V=M_5^2|\Phi|^2+M_D^2|X|^2+\kappa_5|\Phi|^2|X|^2\geq0
\end{equation}
has its minimum at $\Phi=X=0$, and spatial-gradient energies are
nonnegative. There is no tree-level condensate, tadpole or quadratic
mixing to alter G1's free tower. This is a classical stability statement,
not a theorem about a renormalized quantum effective potential.
Dimensions are $[\Phi]=[X]=3/2$ and $[\kappa_5]=-1$.
The interacting theory is a five-dimensional effective field theory
(EFT), not a proposed ultraviolet completion.

\section{The vertex fixes the momentum transfer}
Using the \emph{same} normalized modes for both fields,
\begin{align}
 \Phi(x,\theta)&=\sum_n\phi_n(x)\frac{e^{\ii n\theta}}{\sqrt{2\pi R}},
 &m_n^2&=M_5^2+\frac{(n+\alpha)^2}{R^2},\\
 X(x,\theta)&=\sum_l x_l(x)\frac{e^{\ii l\theta}}{\sqrt{2\pi R}},
 &\mu_l^2&=M_D^2+\frac{l^2}{R^2}.
\end{align}
The four-mode overlap is exact:
\begin{equation}
 \kappa_5\int_0^{2\pi}R\dd\theta\,
 u_m^*u_nu_j^*u_l
 =\frac{\kappa_5}{2\pi R}\delta_{n+l,m+j}
 \equiv g\,\delta_{n+l,m+j}.
\end{equation}
Thus one dimensionless coupling determines all tree vertices,
\begin{equation}\label{eq:vertex}
 \cL_{\rm int}^{(4)}
 =-g\sum_{m,n,j,l}\delta_{n+l,m+j}
       \phi_m^\dagger\phi_n x_j^\dagger x_l,\qquad
 g=\frac{\kappa_5}{2\pi R}.
\end{equation}
No mode-specific conversion matrix is chosen.
With covariant state normalization and
$\langle f|S-1|i\rangle=\ii(2\pi)^4\delta^4(P_f-P_i)\cM$,
the distinct-scalar process is
\begin{equation}\label{eq:process}
 \phi_n+X_l\longrightarrow\phi_m+X_j,\qquad
 j=l+n-m,\qquad \boxed{\cM_{\rm tree}=-g}.
\end{equation}
There is no identical-final-particle factor $1/2$.
The four-dimensional energy and momentum are conserved. Compact
momentum conservation is equally explicit:
\begin{equation}
 \frac{n+\alpha}{R}+\frac lR
 =\frac{m+\alpha}{R}+\frac jR,\qquad
 \Delta p_{5,X}=\frac{n-m}{R}=-\Delta p_{5,\Phi}.
\end{equation}
The canonical translation generator conserves $n+l$; adding the
unchanged charged-particle holonomy shift gives the kinetic-momentum
version above. The detector is the momentum receiver.
Large gauge relabeling
$(\alpha,n,m)\mapsto(\alpha+k,n-k,m-k)$ leaves every formula invariant.

\section{Why heavy recoil opens the previously closed channel}
G1's candidate $\phi_n\to\phi_m+\chi_{n-m}$, with a massless bulk
$\chi$, costs at least $|n-m|/R$ in the emitted mode. For $M_5>0$,
$m_n-m_m<|n-m|/R$, so that decay is closed.
The present process is \emph{scattering}, with an incoming detector.
For its $l=0$ component, write
$d=n-m\ne0$ and $\Delta=m_n-m_m>0$. The detector's rest-energy
increase is instead
\begin{equation}
 \delta E_X=\sqrt{M_D^2+d^2/R^2}-M_D
 =\frac{d^2/R^2}{\sqrt{M_D^2+d^2/R^2}+M_D}.
\end{equation}
It tends to $d^2/(2M_D R^2)$, not $|d|/R$, for a heavy detector.
An exothermic channel at vanishing ordinary incoming kinetic energy
exists precisely when
\begin{align}
 \Delta&>\delta E_X,\\
 \boxed{M_D>\frac{d^2/R^2-\Delta^2}{2\Delta}} .
 \label{eq:critical}
\end{align}
Squaring is legitimate because $M_D+\Delta>0$.
For the G1 massive tower, $0<\Delta<|d|/R$, so the right side is
positive. Equality is threshold only; nonzero incident kinetic energy
may open a channel that fails this rest-limit condition.
This is a finite-recoil mechanism, not an unrecorded supply of energy.

\section{Cross section, thresholds and inverse process}
For a component $l$ with center-of-mass incoming momentum $p_i>0$,
\begin{align}
 E_n&=\sqrt{m_n^2+p_i^2},&
 E_l&=\sqrt{\mu_l^2+p_i^2},&
 \sqrt{s_l}&=E_n+E_l,\\
 p_f&=\frac{\sqrt{[s_l-(m_m+\mu_j)^2]
                         [s_l-(m_m-\mu_j)^2]}}{2\sqrt{s_l}}.
 \label{eq:pf}
\end{align}
Always test $\sqrt{s_l}>m_m+\mu_j$ \emph{before} using the square
root; its product alone can also be positive below the pseudothreshold.
Otherwise the physical open-channel cross section is zero.

For the open channel, integrate the invariant two-body phase space
with flux $4p_i\sqrt{s_l}$:
\begin{align}
 \frac{\dd\sigma_l}{\dd\Omega}
 &=\frac{|\cM|^2}{64\pi^2s_l}\frac{p_f}{p_i},\\
 \boxed{\sigma_l=\frac{g^2}{16\pi s_l}\frac{p_f}{p_i}},\qquad
 v_{{\rm M},l}&=p_i\left(\frac1{E_n}+\frac1{E_l}\right),\\
 K_l\equiv\sigma_l v_{{\rm M},l}
 &=\boxed{\frac{g^2p_f}{16\pi E_n E_l\sqrt{s_l}}}.
 \label{eq:rate}
\end{align}
This is the usual relativistic two-body normalization
\cite{pdg,tong}; the contact amplitude is isotropic at this order.
Both $\sigma_l$ and $K_l$ have dimension $-2$. Neither is a
dimensionless probability.
In an exothermic channel, $\sigma_l$ grows as $1/p_i$ as
$p_i\to0^+$, while $K_l$ has a finite limit. No flux exists at
exactly zero relative velocity for defining $\sigma$ by division.
For the inverse process at the \emph{same} $s_l$,
\begin{equation}
 p_i^2\sigma_{n,l\to m,j}
 =p_f^2\sigma_{m,j\to n,l}
 =\frac{g^2p_i p_f}{16\pi s_l}.
\end{equation}
Preparing the inverse channel at the same small numerical
momentum $0.2$ instead is a different experiment and can be below
threshold; this does not violate detailed balance.

\subsection*{One predeclared engineering example}
Take $R=1$, $M_5=0.5$, $\alpha=0.25$ from the G1 quarter-holonomy
card and add $M_D=5$, $g=0.05$, $p_i=0.2$. These are arbitrary
common energy units, not measured particle data. Consider
$n=1,l=0\to m=0,j=1$:
\begin{center}
\begin{tabular}{ll}
\toprule
Quantity & Value\\
\midrule
$(m_1,m_0)$ & $(1.34629120,\ 0.55901699)$\\
$(\mu_0,\mu_1)$ & $(5,\ 5.09901951)$\\
Critical $M_D$ from \eqref{eq:critical} & $0.24146563$\\
$\sqrt s,\ p_f$ & $6.36506416,\ 1.02175458$\\
$\sigma_{1,0\to0,1}$ & $6.27164081\,10^{-6}$\\
$K_{1,0\to0,1}$ & $1.17224311\,10^{-6}$\\
\bottomrule
\end{tabular}
\end{center}
The energy ledger is
\begin{align}
 m_1-m_0&=0.78727421,& \mu_1-\mu_0&=0.09901951,\\
 T_{\rm out}-T_{\rm in}
 &=0.68825469,& T_{\rm in}&=0.01877296,\\
 T_{\rm out}&=0.70702765.&
\end{align}
Here $T$ is the total \emph{ordinary three-dimensional} kinetic energy
of both particles; internal kinetic energy is already in their
four-dimensional rest masses. At the event,
$1.25+0=0.25+1$ in compact momentum units and
$E_n+E_l=E'_m+E'_j$. The final ordinary kinetic energy receives
the difference. No work reservoir has been omitted.

\section{A localized detector is a state, not an external drive}
Choose the finite, exactly normalized preparation
\begin{equation}\label{eq:packet}
 c_l=\mathcal N e^{-l^2/(4w^2)}e^{-\ii l\theta_0},
 \quad l=-4,\ldots,4,\qquad
 w=1.2,\quad\theta_0=0.7,\quad \sum_l|c_l|^2=1 .
\end{equation}
Finite support defines this particular state; it is not an
uncontrolled replacement of an infinite Gaussian. At a common
narrow ordinary momentum centered on $-\boldsymbol p_i$, its
canonical angular wavefunction is
\begin{equation}
 \psi_{\rm ang}(\theta,t)=\frac1{\sqrt{2\pi R}}
 \sum_l c_l e^{\ii l\theta-\ii E_l t},\qquad
 \int R\dd\theta\,|\psi_{\rm ang}|^2=1.
\end{equation}
It is localized initially and freely spreads; there is no potential
holding it fixed. Its first angular moment is
\begin{equation}
 \langle e^{\ii\theta}\rangle_t
 =\sum_l c_{l+1}^*c_l e^{\ii(E_{l+1}-E_l)t}.
\end{equation}
This norm is a canonical one-particle localization diagnostic, not
an assumed Lorentz-covariant local probability density.
For comparison, the local normally ordered field-density profile
in this narrow-momentum limit is proportional to
\begin{equation}
 \langle :X^\dagger X:(\theta,t)\rangle
 \ \propto\
 \left|\sum_l\frac{c_l}{\sqrt{2E_l}}\,
          \frac{e^{\ii l\theta-\ii E_l t}}{\sqrt{2\pi R}}\right|^2 .
\end{equation}
It is not normalized to unit probability by itself. The heavy-field
limit suppresses the distinction between the two profile shapes.

\subsection*{Exact first-order wavepacket statement}
Let $\dd\Pi_a(p)=\dd^3p/[(2\pi)^3\,2E_a(p)]$.
For normalized four-dimensional envelopes $f(p)$, $h_l(k)$, form
\begin{equation}
 |\mathrm{in}\rangle=\sum_l c_l\int\dd\Pi_n(p)\dd\Pi_l(k)
 f(p)h_l(k)\,|\phi_n,p;X_l,k\rangle ,
 \quad \int\dd\Pi_n|f|^2=\int\dd\Pi_l|h_l|^2=1.
\end{equation}
At first order, the outgoing amplitude for a specified $(m,j)$ is
\begin{equation}\label{eq:waveamp}
 A_{mj}(p',k')=-\ii g\,c_{j-n+m}
 \int\dd\Pi_n(p)\dd\Pi_{j-n+m}(k)\,
 f(p)h_{j-n+m}(k)(2\pi)^4\delta^4(p+k-p'-k').
\end{equation}
The coefficient is zero if its index is outside the preparation.
Nonzero smooth envelopes around the open example shell give
nonzero amplitude; ordinary localization is supplied by these
envelopes, not a background force. A positive envelope choice near
that shell avoids an accidental destructive cancellation.
This proves existence of a genuine packet conversion at leading
order, without treating a plane-wave cross section as a probability.

For conversion $m\ne n$, when detector recoil is unobserved and only
outgoing $m$ is selected,
orthogonality of distinct $j=l+n-m$ gives
\begin{equation}\label{eq:trace}
 P_m^{\rm LO}=\sum_j\int\dd\Pi_m(p')\dd\Pi_j(k')\,|A_{mj}|^2
 =\sum_l |c_l|^2 P_{n,l\to m,l+n-m}^{\rm LO}[f,h_l].
\end{equation}
It is not $|\sum_l c_l\cM_l|^2$. Absolute encounter probabilities
require the envelopes and encounter geometry; no such probability
is claimed from the engineering card alone.
Translating $\theta_0$ changes phases but not \eqref{eq:trace}:
an incident single $n$ mode is uniform around the circle.
A dephased detector with the same $|c_l|^2$ gives the same inclusive
conversion; that observable alone cannot certify localization.
Coherent incoming $\Phi$ modes or coherent recoil post-selection
would define a different experiment.

\subsection*{The precisely scoped numerical packet rate}
The kinetic approximation uses component overlap integrals
$I_l=\int\dd t\,\dd^3x\,n_\Phi(x,t)n_{X,l}(x,t)$ with
normalized conditional spatial profiles. It gives
\begin{equation}
 N_m^{\rm LO}\simeq\sum_l |c_l|^2 I_l K_{n,l\to m,l+n-m}.
\end{equation}
Only under a common-overlap approximation $I_l=I$ can one write
\begin{equation}
 N_m^{\rm LO}\simeq I K_{\rm packet}^{(m)},\qquad
 K_{\rm packet}^{(m)}=\sum_l |c_l|^2 K_{n,l\to m,l+n-m}.
\end{equation}
Each component has its own $s_l$ and velocity. The recorded
$K_{\rm packet}$ is this rate coefficient, not a universal
wavepacket cross section or a probability. The approximation need
not hold during substantial spreading; \eqref{eq:waveamp} is the
appropriate starting point then.
All open final $m$ are enumerated before defining a target rate
fraction: a sufficient finite search interval is
$|m+\alpha|\leq R\sqrt{s_l}$, followed by the exact threshold test.
The elastic $m=n$ channel is included in the total collision rate.
Its scattering cross section is not the total probability of staying
in the original KK sector: that probability also contains the identity
part of $S$ and its interference, omitted from the connected amplitude
\eqref{eq:waveamp}. No G1 probe weight is reused as a scattering fraction.

For the stated preparation the common-overlap coefficients are
$K_{\rm packet}^{(0)}=1.11933296\,10^{-6}$ and
$K_{\rm packet}^{\rm all}=1.85661445\,10^{-6}$.
Their ratio $0.60288929$ is a target fraction of \emph{scattering
rates under this prescription}, not a universal conversion probability.
At $l=0$ the open final branches are $m=-1,0,1$.

\section{Verification, limits and next physical decision}
The companion verifier checks circle overlap, independent energy-shell
roots and phase space, thresholds, conservation laws, inverse processes
and gauge relabeling. It also checks packet spreading, trace/dephasing
identities and full open-channel enumeration. Reports preserve the
source hash and every assertion. G2 passes 262/262 implementation
checks on independent rerun; G1 passes its separate 331/331 regression.
The finite packet has $|\langle e^{\ii\theta}\rangle|=0.91640536$
at $t=0$ and $0.48041872$ at $t=5$, with conserved norm and
mean energy. This verifies a freely evolving localized preparation,
not a localized vacuum.

The engineering amplitude is tree level. Weak $g$ is not by itself
a bound on the error of a five-dimensional EFT: a physical use
requires a cutoff above all occupied and accessible momenta,
matched higher operators, and loop/counterterm control. At
nonzero coupling the exact interacting propagator has self-energy
corrections and cuts. G1's free pole and probe formulae are the
quadratic baseline, not all-orders identities of \eqref{eq:action}.
No numerical EFT-error bar, radiative naturalness, background
stabilization, Standard Model assignment or Route-F fit is certified.

\paragraph{Physical result.}
An incoming excitation can leave in another existing mass branch
through a local collision. The detector takes the compensating
compact momentum and ordinary recoil. Neither the shared bulk
mass $M_5$ nor a coordinate convention changes.
The simple G1 decay obstruction is bypassed by changing the
\emph{physical process}, not violating its inequality.

\paragraph{Bounded next options.}
The most direct continuation is a recoil-resolved experiment:
test the correlation $j-l=n-m$ together with the predicted ordinary
outgoing momentum. A finite-encounter calculation can specify
four-dimensional packets and evaluate \eqref{eq:waveamp} if an
absolute probability is needed. A driven background is a backup
only, with work included in the Hamiltonian ledger. Stabilization
and particle-content questions remain separate G3/G4 hypotheses;
they do not reopen this conditional tree-level calculation.

\begin{thebibliography}{9}
\bibitem{pdg}
Particle Data Group, \emph{Kinematics}, Review of Particle Physics
(2024), sections 49.5--49.5.1, equations (49.27)--(49.33),
\href{https://pdg.lbl.gov/2024/reviews/rpp2024-rev-kinematics.pdf}{official review}.
\bibitem{tong}
D.~Tong, \emph{Quantum Field Theory}, sections 3.4 and 3.6,
\href{https://www.damtp.cam.ac.uk/user/tong/qft/qfthtml/S3.html}{author's lecture notes}.
\end{thebibliography}
\end{document}
